\section{What is good：模型好坏有多个定义}

这一节把“模型好”拆成几种互相竞争的定义。一个模型可以 benchmark 强但贵，一个模型可以便宜但不被用户喜欢，一个模型可以在 Arena 上受欢迎但在专业任务上错误很多。老师用这组例子说明：评价本身就是价值选择。

接下来四张图不是并列展示四个网站，而是在构造一条论证链：先看“标准化能力”，再看“成本约束”，再看“用户偏好”，最后看“市场选择”。如果这四个排序一致，评价会很简单；现实中它们经常不一致，所以我们才需要明确评价目的。

\begin{figure}[H]
\centering
\includegraphics[width=0.9\textwidth]{artificial-analysis.png}
\caption{Artificial Analysis leaderboard：按 benchmark 和能力指标比较模型。}
\figsource{CS336 Lecture 12 官方源引用图，本地化自 Stanford lectures images。}
\end{figure}

\begin{knowledgebox}{读图：benchmark leaderboard 测的是“标准化能力”}
这类 leaderboard 让模型在同一批标准任务上比较，优点是清晰、可复现、便于购买决策；缺点是指标集合会选择性代表“好”，且容易被过拟合、污染或专门优化。读这类图时要先问：它覆盖了哪些任务，没覆盖哪些真实场景。图中的排序不是模型的绝对价值，而是某组 benchmark 权重下的投影。
\end{knowledgebox}

这个图回答的是“如果我只关心一组标准能力，谁强？”但真实部署还会问另一个问题：强多少，贵多少，慢多少？因此老师接着把 cost 引入评价。

\begin{figure}[H]
\centering
\includegraphics[width=0.9\textwidth]{artificial-analysis-cost.png}
\caption{加入成本后的模型比较：好模型不只要强，还要便宜、快、可部署。}
\figsource{CS336 Lecture 12 官方源引用图。}
\end{figure}

\begin{knowledgebox}{读图：性能-成本 Pareto frontier}
如果两个模型能力相近，成本更低的模型可能更“好”。如果一个模型略强但贵很多，具体是否值得取决于用户任务价值。读这张图时应该找 Pareto frontier：有没有模型在更便宜的同时不牺牲太多能力？有没有模型虽然榜单高，但价格或 latency 让它不适合实际产品？
\end{knowledgebox}

第三种定义来自用户偏好。标准 benchmark 假定题目和答案都由评测方定义；Arena 则让真实用户决定两个回答哪个更好。这更接近开放式 assistant 使用，但也更混乱。

\begin{figure}[H]
\centering
\includegraphics[width=0.58\textwidth]{lmarena-leaderboard.png}
\caption{LM Arena leaderboard：通过人类 pairwise preference 给模型排序。}
\figsource{Arena AI / Chatbot Arena leaderboard，本地化后供 CS336 Lecture 12 使用。}
\end{figure}

\begin{knowledgebox}{读图：preference leaderboard 测的是“人更喜欢哪个回答”}
人类偏好能捕捉开放式回答的风格、完整性和实用性，但也混合了正确性、语气、长度、迎合性和用户群体偏见。它不是纯粹能力测量，而是特定交互分布上的偏好测量。老师在这里的提醒是：真实用户信号很宝贵，但它不是干净标签。
\end{knowledgebox}

最后一种定义更商业：用户是否真的选择并付费使用。OpenRouter 这类 ranking 把模型放进市场选择中，但市场选择也受价格、可得性、默认选项和生态影响。

\begin{figure}[H]
\centering
\includegraphics[width=0.78\textwidth]{openrouter.png}
\caption{OpenRouter ranking：用户选择和付费也构成一种模型评价信号。}
\figsource{OpenRouter ranking，本地化后供 CS336 Lecture 12 使用。}
\end{figure}

\begin{warningbox}{读图：使用量不是能力本身}
用户选择受到价格、品牌、可用性、上下文长度、API 兼容性、区域访问、产品默认值等影响。商业使用量能说明真实 adoption，但不能直接等同于模型能力。一个模型可能因为便宜、稳定或集成方便而使用量高，不一定因为它在所有任务上最强。
\end{warningbox}

\subsection{本章小结}

模型“好”可以指 benchmark 强、单位成本低、用户偏好高、市场使用多。这些定义都合理，但回答的问题不同。下一节转向最传统的 LM 指标 perplexity，它曾经是语言模型评价的中心，但会暴露“指标平滑”和“真实使用”之间的张力。

